\chapter{Manipulação de Dados}

\section{\texttt{generate} e \texttt{mutate}}

Criam ou substituem colunas. As duas formas são equivalentes em resultado; diferem
na semântica de mutação:

\begin{center}
\begin{tabular}{lll}
\toprule
\textbf{Forma} & \textbf{Semântica} & \textbf{Exemplo} \\
\midrule
\texttt{generate df col = expr}     & Imperativa (muta \texttt{df}) & estilo Stata \\
\texttt{let df2 = generate(df, col = expr)} & Funcional (df inalterado) & estilo funcional \\
\texttt{df |> generate(col = expr)} & Pipe assign-back (muta \texttt{df}) & encadeamento \\
\bottomrule
\end{tabular}
\end{center}

\begin{lstlisting}[caption={generate: forma imperativa}]
load "dados.csv" as df

generate df lrenda = log(renda)
generate df renda2 = renda^2
generate df dummy = if(renda > 1000, 1, 0)
\end{lstlisting}

\begin{lstlisting}[caption={mutate: forma funcional}]
let df2 = mutate(df, lrenda = log(renda), renda2 = renda^2)
// df não possui lrenda nem renda2
// df2 as possui
\end{lstlisting}

\section{Geradores de variáveis aleatórias}

Todas as distribuições abaixo estão disponíveis dentro de expressões
\lstinline{generate}. Elas respeitam a semente global definida por
\lstinline{set_seed} e retornam um valor por linha, propagando parâmetros
escalares para o comprimento do DataFrame.

\begin{lstlisting}[caption={Simulando dados}]
set_seed(42)
let df = dataframe(1000)

generate df u = rnormal(0, 1)
generate df treated = rbernoulli(0.5)
generate df duration = rexponential(1)
\end{lstlisting}

\textbf{Contínuas:} \lstinline{rnormal}, \lstinline{rlognormal}, \lstinline{rskewnormal}, \lstinline{rcauchy}, \lstinline{rstudentt} (alias \lstinline{rt}), \lstinline{rchisq}, \lstinline{rf}, \lstinline{rbeta}, \lstinline{rgamma}, \lstinline{rexponential}, \lstinline{rweibull}, \lstinline{rpareto}, \lstinline{rpert}, \lstinline{rtriangular}, \lstinline{rfrechet}, \lstinline{rgumbel}, \lstinline{rinversegaussian}, \lstinline{rnig}, \lstinline{runiform}.

\textbf{Discretas:} \lstinline{rbernoulli}, \lstinline{rbinomial}, \lstinline{rpoisson}, \lstinline{rgeometric}, \lstinline{rhypergeometric}, \lstinline{rzeta}, \lstinline{rzipf}.

\begin{nota}
  \lstinline{rgeometric(p)} retorna o número de falhas antes do primeiro
  sucesso, com suporte $\{0, 1, 2, \ldots\}$.
\end{nota}

\section{\texttt{filter}}

Seleciona linhas que satisfazem uma condição:

\begin{lstlisting}
let recentes = filter(df, ano >= 2010)
let ricos = filter(df, renda > 5000 and uf == "SP")
\end{lstlisting}

A forma imperativa usa \lstinline{keep}:

\begin{lstlisting}
keep df if ano >= 2010
\end{lstlisting}

\section{\texttt{select} e \texttt{drop}}

\begin{lstlisting}
// mantém só essas colunas
let slim = select(df, ano, uf, renda)
let sem_id = drop(df, id, codigo)        // remove essas colunas
\end{lstlisting}

\section{\texttt{rename}}

\begin{lstlisting}
rename df antiga nova
// renomeia renda_familiar → renda
rename df renda_familiar renda
\end{lstlisting}

\section{\texttt{sort}}

\begin{lstlisting}
sort df ano               // ascendente
sort df renda desc        // descendente
sort df uf ano desc       // múltiplas colunas
\end{lstlisting}

\section{\texttt{dropna}}

Remove linhas com valores ausentes:

\begin{lstlisting}
dropna df                 // remove linhas com qualquer NaN
// remove linhas com NaN em renda ou gini
dropna df renda gini
\end{lstlisting}

\section{\texttt{ffill}}

Forward-fill substitui \texttt{NaN} pelo último valor não ausente. Útil quando
uma série mensal ou trimestral precisa ser expandida para frequência diária:

\begin{lstlisting}[caption={Preenchendo uma taxa mensal em dados diários}]
let daily = merge(asset, rf, key=date, type=left)
let daily = ffill(daily)
\end{lstlisting}

\section{\texttt{append}}

Empilha dois DataFrames verticalmente:

\begin{lstlisting}
let df_total = append(df_2024, df_2025)
\end{lstlisting}

As colunas devem ser compatíveis. Colunas ausentes em um dos DataFrames são
preenchidas com \texttt{NaN}.

\section{\texttt{rbind}}

Concatena uma \textit{lista} de DataFrames verticalmente em uma única
passada. Entradas \texttt{nil} são silenciosamente ignoradas --- útil ao
combinar resultados de \texttt{parallel for} onde algumas iterações
retornam \texttt{nil}:

\begin{lstlisting}[caption={rbind com parallel for}]
let resultados = parallel for i in 0..n, threads=8 {
  let t = tickers[i]
  if t == "SPY" { return nil }
  let df_t = calcular_betas(t)
  df_t
}
let todos_betas = rbind(resultados)
\end{lstlisting}

\section{\texttt{collapse} e \texttt{group\_by}}

\lstinline{group_by} agrega por grupos:

\begin{lstlisting}[caption={Agregação por grupo}]
let media_uf = group_by(df, uf,
  renda_media = mean(renda),
  n           = count()
)
\end{lstlisting}

\lstinline{collapse} é a forma imperativa estilo Stata:

\begin{lstlisting}
collapse df (mean) renda gini, by(uf ano)
\end{lstlisting}

\section{\texttt{pivot\_longer} e \texttt{pivot\_wider}}

\begin{lstlisting}[caption={Formato largo para longo}]
let longo = pivot_longer(df,
  cols  = ["t1", "t2", "t3"],
  names = "trimestre",
  values = "valor"
)
\end{lstlisting}

\begin{lstlisting}[caption={Formato longo para largo}]
let largo = pivot_wider(df,
  id     = "municipio",
  names  = "ano",
  values = "pib"
)
\end{lstlisting}

\section{\texttt{replace}}

Substitui valores condicionalmente:

\begin{lstlisting}
// zera valores negativos
replace df renda = 0 if renda < 0
replace df gini = . if gini > 1         // marca como missing
\end{lstlisting}

\section{\texttt{tabulate}}

Tabela de frequências:

\begin{lstlisting}
tabulate df uf
tabulate df uf ano    // tabulação cruzada
\end{lstlisting}

\section{Auxiliares de arquivo}

Pequenos auxiliares para scripts que persistem dados ou verificam cache:

\begin{lstlisting}[caption={Auxiliares de cache}]
ensure_dir("cache")                    // cria diretório se não existir
if file_exists("cache/dados.csv") {
    load "cache/dados.csv" as df
} else {
    let df = download(...)
    export(df, "csv", "cache/dados.csv")
}

write("texto bruto", "nota.txt")     // escreve string em arquivo
\end{lstlisting}

\section{\texttt{rolling}}

\texttt{rolling} estima OLS em uma janela móvel. É útil para visualizar
coeficientes que variam no tempo ou testar a estabilidade de parâmetros.
Uma coluna \texttt{date} opcional armazena o rótulo do final de cada
janela.

\begin{lstlisting}[caption={OLS rolante com rótulos de data}]
let roll = rolling(y ~ x1 + x2, df, window=252, date=date)
let betas = tidy(roll)
print(betas)
\end{lstlisting}

\section{\texttt{tidy} e \texttt{glance}}

Essas funções convertem objetos de modelo em DataFrames. \texttt{tidy}
retorna um DataFrame organizado com coeficientes, erros-padrão e
intervalos de confiança. Para modelos \texttt{rolling}, retorna uma
linha por janela com colunas \texttt{date}, \texttt{r2} e uma coluna de
coeficiente por variável. \texttt{glance} retorna um resumo de uma
linha com estatísticas de ajuste.

Ambas as funções suportam todos os tipos de modelo: OLS, IV, logit,
probit, ordenado, multinomial, painel (FE, RE, PCSE, GLS, GLSAR,
Arellano-Bond, SystemGMM), séries temporais (ARIMA, AutoReg, ARDL,
GARCH, VAR, VECM, VARMA, ETS, Kalman), sobrevivência (Cox, KM), causal
(DiD, Threshold), multivariada (SUR, 3SLS), regularização (ridge, lasso,
elasticnet), entre outros.

\begin{lstlisting}[caption={Tidy e glance}]
let m = ols(y ~ x1 + x2, df)
let coefs = tidy(m)
let summary = glance(m)
\end{lstlisting}
